%---------------------------------------------------------------
% SLO: slovenski povzetek
% ENG: slovenian abstract
%---------------------------------------------------------------
\selectlanguage{slovene} % Preklopi na slovenski jezik
\addcontentsline{toc}{chapter}{Povzetek}
\chapter*{Povzetek}

\noindent\textbf{Naslov:} \ttitle
\bigskip

Magistrsko delo primerja neposredno sklepanje iz konteksta in pretvorbo besedila v SQL pri pridobivanju
informacij iz strukturiranih podatkov z velikimi jezikovnimi modeli. Ovrednotili smo \num{10} lokalno
nameščenih modelov z odprtimi utežmi in \num{6} serializacijskih formatov na podatkovnih bazah iz zbirke
Spider in podvzorčenih bazah iz zbirke BIRD. Če neuspele zagone in zagone s preseženim kontekstnim oknom
ovrednotimo z ničlo, je pretvorba besedila v SQL najboljši neposredni format v povprečnem normaliziranem
natančnem ujemanju v celoti na kohorti odgovorljivih vprašanj presegla za \num{46.5} oziroma \num{44.7}
odstotne točke, na preseku s kontekstno ustrezno kohorto, kjer vse konfiguracije ustrezajo kontekstnemu
oknu, pa za \num{45.8} oziroma \num{38.0} odstotne točke; pri tem je bila ustrezna prednost po meri
BERTScore na tem preseku le \num{4.9} oziroma \num{2.0} odstotne točke. Primerjave operativne učinkovitosti
na kontekstno ustrezni kohorti so pri pretvorbi besedila v SQL pokazale manjšo porabo vhodnih žetonov in
krajši povprečni odzivni čas. Med neposrednimi formati je bil najučinkovitejši CSV, sledil mu je TOON
(Token-Oriented Object Notation). Ugotovitve so podlaga za odločitveni okvir, ki daje prednost izbiri
metode pred serializacijo ter priporoča pretvorbo besedila v SQL, kadar je izvajanje poizvedb SQL na voljo
in primerno za nalogo.

\subsection*{Ključne besede}
\textit{\tkeywords}
\clearemptydoublepage
